\chapter{Discussion}
\label{ch:discussion}

\section{What the strict split reveals}
\label{sec:disc-strict}

The central methodological result is that a disclaimer became a measurement.
Most audio classification on found data either ignores speaker and channel
leakage or mentions it qualitatively. Here the pseudo-speaker split bounds it:
12--14 macro-F1 points for pretrained representations and about 25 for
hand-crafted features. Three implications follow.

\begin{enumerate}
  \item \textbf{Scores reported on clip-level splits of found-data corpora are
        inflated, and by an amount that depends on the representation.}
        Comparing papers that use different features on clip-level splits
        partly ranks how well each exploits the channel, not how well it
        identifies dialect.
  \item \textbf{The honest acoustic headline for this corpus is 0.808}
        (fine-tuned XLS-R, strict split). That remains a strong 9-way result,
        about seven times chance, on speakers the model has never heard.
  \item \textbf{Per-class strict numbers must be read with care.} Barishal's
        strict test set is a single 251-clip pseudo-speaker cluster, so its fall
        from 0.871 to 0.603 for the frozen model measures generalisation to one
        speaker or source rather than to the class. Conclusions at the macro
        level are the robust ones.
\end{enumerate}

\section{Errors follow dialect geography}
\label{sec:disc-geography}

Under every model and protocol the dominant confusion is between Barishal and
Noakhali, two southern coastal dialects with documented linguistic proximity,
and Mymensingh, a transitional zone, spreads its errors across its neighbours.
This is linguistically coherent behaviour: the models' mistakes recapitulate the
dialect continuum rather than arbitrary decision boundaries, which supports the
claim that they learn accentual structure and not only channel shortcuts, even
on the clip-level split.

\section{The shrinking advantage of fine-tuning}
\label{sec:disc-finetune}

Fine-tuned XLS-R gains 3.5 points over frozen Whisper at clip level, a
significant difference, but only 1.5 under the strict split, which is not
significant, and its strict-split drop ($-14.4$) is slightly larger than the
frozen model's ($-12.4$). End-to-end adaptation lets a model exploit whatever
separates the training classes, including channel signatures that correlate with
class in the training data; a frozen, general-purpose representation cannot
specialise in this way. The practical lesson is that, when leakage cannot be
ruled out, frozen-embedding baselines are not merely cheaper: they partly
protect against shortcut learning, and a fine-tuning margin measured on a leaky
split overstates the improvement that would hold in deployment. The comparison
is confounded, however: the frozen model is Whisper-small and the fine-tuned one
XLS-R-300M, so ``frozen versus fine-tuned'' and ``Whisper versus XLS-R'' are not
separated (\cref{sec:disc-limits}).

\section{The augmentation result: prosody is signal}
\label{sec:disc-augment}

The channel-suppression chain lowered strict-split macro-F1 by six points. The
chain included pitch shift ($\pm 2$ semitones) and time stretch (0.9--1.1),
transforms that are standard and beneficial in speech recognition
\citep{ko2015augmentation}, where the transcript is invariant to prosody.
Dialect identification inverts this: pitch contour, rhythm and tempo carry
dialect, so these transforms corrupt the evidence for the label rather than
removing nuisance variation. Probing work showing that accent fine-tuning
reshapes prosodic encoding in wav2vec~2.0 \citep{yang2023accent} supports this
reading, and the validation score, depressed throughout training (0.702 against
0.757), indicates a harder target rather than an unlucky checkpoint.

The recommendation is task-specific: augmentation for dialect and accent
identification should preserve prosody (additive noise, equalisation, band
limiting, gain), and prosodic perturbation should be treated as harmful until
shown otherwise. The reduced chain was not tested and is left as future work;
without that test, the attribution to the pitch and tempo transforms remains a
hypothesis.

\section{Words or sound?}
\label{sec:disc-words}

The lexical route reads the thesis title literally: how much of Bangla dialect
identity is in the words, and how much lies beyond them? Three findings stand
out.

\paragraph{Standard-Bangla speech recognition removes the dialect before a
classifier can see it.} Recognisers trained on standard Bangla (Tugstugi,
Hishab) transcribe regional speech into standard forms---\bn{biya}{বিয়া দিয়া দাও}
(``arrange the marriage'') becomes \bn{biye}{বিয়ে দিয়ে দাও}---and reproduce only
1--3\,\% of the dialect words in Ben-10's human transcripts. Their output on
regional clips is barely less standard than on Formal clips. The one front-end
fine-tuned on dialect speech reproduces 19\,\%, which is both the best available
and a hard ceiling for any word-list method: four of five dialect words never
reach the text. This confirms, with an automated measure over the full
validation set, the observation that \citet{dipto2025ben10} made by hand, and it
is the single most important fact about the lexical route.

\paragraph{Filtering to regional words discards evidence.} The intuitive
design---drop every word that standard Bangla also has and match what is left
against regional word lists---scores 0.360 under the strict split, because
fewer than half of the 5-second clips keep any matchable word. A TF-IDF
classifier over all words and character n-grams reaches 0.633 on the same
transcripts. Two mechanisms plausibly explain the gap, though this thesis did
not isolate them: words that also exist in standard Bangla but are used at
different rates across regions (pronouns, particles, common verbs), which the
filter removes; and character n-grams, which generalise over rare inflected
dialect forms such as verb endings that never reach the lexicon's frequency
thresholds. Word matching keeps its value as an \emph{explanation}---it shows
which words decided a clip and how strongly each belongs to each region---rather
than as the classifier.

\paragraph{Words and sound are complementary, modestly.} Fusion raises the
honest headline from 0.808 to 0.856, a gain that is borderline under the
speaker-group bootstrap (a 4\,\% probability of no gain), and the clip-level
score from 0.952 to 0.967, which is significant. The gain lands where the
acoustic model is weakest under the strict split---Barishal (+21 points) and
Sylhet (+8), the dialect with the most distinctive and best-transcribed
vocabulary---and is negative only for Khulna, where the acoustics are already
near perfect. The two routes fail on different clips, which is the condition
under which fusion helps, and this is evidence that the text carries dialect
information of its own rather than a noisier copy of the acoustic decision.

\paragraph{The answer.} For clips this short, dialect identity lies
\emph{mostly beyond the words}. Words alone (0.633) sit between hand-crafted
acoustics (0.480) and frozen Whisper (0.793); the acoustic route is the stronger
single route, and the lexical route's ceiling is set by speech recognition, not
by the classifier. Two caveats temper the comparison. The TF-IDF model may
partly learn programme- or topic-specific vocabulary, because the strict split
separates speakers but not content, so 0.633 is itself an upper bound on the
dialect information in the words. And the acoustic models were trained on 13
hours of in-domain audio, whereas the recognisers saw none of it: a
dialect-aware recogniser trained on this corpus would raise the lexical ceiling,
at the cost of the off-the-shelf setting that makes the comparison clean.

\section{Limitations}
\label{sec:disc-limits}

\begin{itemize}
  \item \textbf{Pseudo-speakers are approximate.} Calibration against Formal's
        true sources reaches an ARI of 0.347, and clustering both over-merges (a
        251-clip Barishal cluster) and over-splits (many singletons). The strict
        split removes much, but not provably all, speaker and channel overlap;
        true speaker labels would tighten the estimate in both directions.
  \item \textbf{Uneven strict test cells.} Group-disjoint splitting with
        unequal clusters concentrates some classes' test sets in a few clusters,
        so per-class strict numbers have high variance.
  \item \textbf{Short clips.} Five seconds excludes longer-range dialect cues
        such as discourse markers and phrase-level intonation.
  \item \textbf{Single domain.} The corpus is broadcast and online media;
        generalisation to conversational or telephone speech is untested.
  \item \textbf{Model family is confounded with adaptation regime.} The frozen
        model is Whisper-small and the fine-tuned one XLS-R-300M. Frozen XLS-R
        embeddings, cheap to compute locally, would separate the two factors.
  \item \textbf{Single-seed fine-tuning.} Each XLS-R result is one seed; the
        bootstrap intervals cover test-set sampling noise, not training-seed
        variance.
  \item \textbf{Whisper pooling includes padding.} Whisper's encoder always
        processes 30\,s of input, and the baseline mean-pools all 1,500 frames,
        of which only about 250 cover the 5\,s of speech. The padding
        contribution is the same for every clip, but pooling only speech frames
        would be the cleaner design.
  \item \textbf{Authenticity of found data.} Clips from films and dramas may
        include actors performing a dialect rather than native speakers.
  \item \textbf{No reference transcripts for the corpus.} Recognition quality is
        measured on Ben-10, a different corpus with longer, spontaneous clips
        and only three districts in common; word error rate on this corpus is
        unknown. Hand-transcribing a subset of about 20 clips per class would
        close this gap.
  \item \textbf{Content leakage is not controlled.} The strict split is
        speaker-disjoint, not topic-disjoint; recurring programmes, names and
        topics can appear on both sides and inflate the text classifier. The
        in-domain lexicon's two-speaker rule and the external lexicon mitigate,
        but do not remove, this risk.
  \item \textbf{Front-end substitutions.} The benchmark's ``Wav2Vec2 (SCB)''
        checkpoint is not public, and the wav2vec~2.0 model used instead is a
        different model; Google's recogniser and the benchmark's Ben-10-tuned
        wav2vec~2.0 were not run. The Ben-10-tuned recogniser was trained on ten
        districts, of which only Barishal, Chattogram and Sylhet are classes
        here, and Khulna, Rajshahi and Dhaka have no clean external word list.
  \item \textbf{Validation reuse.} The XLS-R checkpoint was selected on the same
        validation split later used to tune the fusion weight $\alpha$, so
        XLS-R's validation probabilities are mildly optimistic and $\alpha$ may
        lean towards the acoustic model. Test numbers are unaffected, and all
        tuning used validation data only.
  \item \textbf{Greedy, single decoding.} Each front-end's transcripts are one
        greedy decode; beam search or temperature fallback might reduce
        Whisper-large-v3's repetition loops.
\end{itemize}

\section{Future work}
\label{sec:disc-future}

\begin{itemize}
  \item A prosody-preserving augmentation chain on the strict split, the direct
        test of \cref{sec:disc-augment}.
  \item True speaker annotation for a subset, to validate the pseudo-speaker
        estimate.
  \item Larger multilingual encoders (XLS-R-1B, MMS) within the same free-GPU
        budget, and frozen XLS-R embeddings to remove the model-family confound.
  \item Longer inputs, and aggregation of clip decisions to speaker or
        recording level.
  \item Cross-domain evaluation on conversational and telephone speech.
  \item For the lexical route: a dialect-aware recogniser fine-tuned on this
        corpus, since transcription rather than classification is the
        bottleneck; a hand-transcribed subset for word error rate on this
        corpus; a topic-disjoint split to bound content leakage; and joint
        rather than late fusion, feeding text features into the acoustic model.
\end{itemize}
